\uuid{3fL9}
\titre{Optimisation avec contrainte }

\niveau{L2}
\module{ Fonctions de plusieurs variables }
\chapitre{Optimisation}
\sousChapitre{Optimisation sous contrainte}
\theme{Etude de fonctions}
\auteur{Quentin Liard}
\datecreate{ 2026-04-03}
\organisation{AMSCC}
\difficulte{3}

\contenu{

\texte{ 
On considère la fonction $f:\mathbb{R}^2\to\mathbb{R}$ définie par
\[
f(x,y)=x^2+y^2-2x-4y
\]
sur le domaine
\[
D=\{(x,y)\in\mathbb{R}^2\mid x\ge 0,\ y\ge 0,\ x+y\le 2\}.
\]
}

\begin{enumerate}

\item   \question{Justifier que $f$ admet un minimum global et un maximum global sur $D$.}
\indication{}
\reponse{
La fonction $f$ est polynomiale, donc continue sur $\mathbb{R}^2$, en particulier sur $D$.

Le domaine $D$ est fermé et borné dans $\mathbb{R}^2$, donc compact.

Par le théorème de Weierstrass, $f$ admet un minimum global et un maximum global sur $D$.
}

\item   \question{Déterminer les points critiques éventuels de $f$ à l’intérieur de $D$.}
\indication{}
\reponse{
On calcule
\[
\frac{\partial f}{\partial x}(x,y)=2x-2,
\qquad
\frac{\partial f}{\partial y}(x,y)=2y-4.
\]
Un point critique intérieur vérifie
\[
2x-2=0,\qquad 2y-4=0,
\]
soit
\[
x=1,\qquad y=2.
\]
Mais ce point n’appartient pas à $D$ car
\[
1+2=3>2.
\]
Il n’y a donc aucun point critique à l’intérieur de $D$.
}

\item   \question{Étudier $f$ sur le bord défini par $x=0$, avec $0\le y\le 2$.}
\indication{}
\reponse{
Sur ce bord,
\[
f(0,y)=y^2-4y.
\]
On pose
\[
\varphi_1(y)=y^2-4y,\qquad y\in[0,2].
\]
Alors
\[
\varphi_1'(y)=2y-4.
\]
Sur $[0,2]$, on a $\varphi_1'(y)\le 0$, donc $\varphi_1$ est décroissante.
Ainsi
\[
\varphi_1(0)=0,\qquad \varphi_1(2)=-4.
\]
Le maximum sur ce bord vaut $0$ en $(0,0)$, et le minimum vaut $-4$ en $(0,2)$.
}

\item   \question{Étudier $f$ sur le bord défini par $y=0$, avec $0\le x\le 2$.}
\indication{}
\reponse{
Sur ce bord,
\[
f(x,0)=x^2-2x.
\]
On pose
\[
\varphi_2(x)=x^2-2x,\qquad x\in[0,2].
\]
Alors
\[
\varphi_2'(x)=2x-2.
\]
Donc $\varphi_2$ admet un minimum en
\[
x=1,
\]
de valeur
\[
\varphi_2(1)=-1.
\]
Aux extrémités :
\[
\varphi_2(0)=0,\qquad \varphi_2(2)=0.
\]
Ainsi, sur ce bord, le minimum vaut $-1$ en $(1,0)$ et le maximum vaut $0$ aux points $(0,0)$ et $(2,0)$.
}

\item   \question{Étudier $f$ sur le bord défini par $x+y=2$, avec $x\in[0,2]$.}
\indication{}
\reponse{
On pose
\[
y=2-x,\qquad x\in[0,2].
\]
Alors
\[
f(x,2-x)=x^2+(2-x)^2-2x-4(2-x).
\]
Développons :
\[
f(x,2-x)=x^2+4-4x+x^2-2x-8+4x
=2x^2-2x-4.
\]
On pose
\[
\varphi_3(x)=2x^2-2x-4,\qquad x\in[0,2].
\]
Alors
\[
\varphi_3'(x)=4x-2.
\]
Le point critique est
\[
x=\frac12,
\qquad y=\frac32.
\]
On obtient
\[
\varphi_3\left(\frac12\right)=2\cdot\frac14-2\cdot\frac12-4=-\frac92.
\]
Aux extrémités :
\[
\varphi_3(0)=-4,\qquad \varphi_3(2)=0.
\]
Ainsi, sur ce bord, le minimum vaut
\[
-\frac92
\]
au point
\[
\left(\frac12,\frac32\right),
\]
et le maximum vaut $0$ au point $(2,0)$.
}

\item   \question{Conclure en déterminant le minimum global et le maximum global de $f$ sur $D$.}
\indication{}
\reponse{
Comme il n’y a aucun point critique à l’intérieur de $D$, les extrema globaux sont sur le bord.

Récapitulons :
\begin{itemize}
\item sur $x=0$ : min $=-4$, max $=0$ ;
\item sur $y=0$ : min $=-1$, max $=0$ ;
\item sur $x+y=2$ : min $=-\dfrac92$, max $=0$.
\end{itemize}

Donc le minimum global de $f$ sur $D$ est
\[
-\frac92
\]
atteint en
\[
\left(\frac12,\frac32\right),
\]
et le maximum global de $f$ sur $D$ est
\[
0,
\]
atteint aux points
\[
(0,0)\quad\text{et}\quad(2,0).
\]
}

\end{enumerate}

}